\documentclass{article}



% Language setting

% Replace `english' with e.g. `spanish' to change the document language

\usepackage[english]{babel}



% Set page size and margins

% Replace `letterpaper' with `a4paper' for UK/EU standard size

\usepackage[letterpaper,top=2cm,bottom=2cm,left=3cm,right=3cm,marginparwidth=1.75cm]{geometry}



% Useful packages

\usepackage{amsmath}

\usepackage{amssymb}

\usepackage{graphicx}

\usepackage[colorlinks=true, allcolors=black]{hyperref}

\usepackage{titlesec}
\usepackage{xparse}



\NewDocumentCommand{\qcq}{m o o m}{%
	\ensuremath{\forall #1 \IfValueT{#2}{, #2} \IfValueT{#3}{, #3}  \in  \mathbb{#4}}%
}
\NewDocumentCommand{\ext}{m o o m}{%
	\ensuremath{\exists #1 \IfValueT{#2}{, #2} \IfValueT{#3}{, #3}  \in  \mathbb{#4}}%
}



\title{\textbf{Continuité}}

\author{\textbf{Yassin Akkab}}



\newcommand{\R}{\mathbb{R}}
\newcommand{\N}{\mathbb{N}}
\newcommand{\Z}{\mathbb{Z}}
\newcommand{\C}{\mathbb{C}}
\newcommand{\Q}{\mathbb{Q}}
\newcommand{\D}{\mathbb{D}}





\graphicspath{ {./Screenshots/} }




\begin{document}
	
	\maketitle
	
	
	
	\tableofcontents
	
	\newpage
	
	\section{Introduction aux fonctions exponentielles}
	
	\subsection{Définition intuitive}
	
	La fonction exponentielle modélise de nombreux phénomènes de croissance ou de décroissance. Notée \( e^x \), cette fonction est définie pour tout réel \( x \), où \( e \) est un nombre particulier, appelé base des logarithmes naturels, et vaut environ \( 2.718 \).
	
	\subsection{Définition formelle}
	
	Formellement, la fonction exponentielle est la fonction qui satisfait l'équation différentielle suivante :
	\[
	y' = y \quad \text{avec} \quad y(0) = 1.
	\]
	Autrement dit, la fonction exponentielle est la seule fonction qui est égale à sa propre dérivée.
	
	\subsection{Propriétés fondamentales}
	
	La fonction exponentielle présente plusieurs propriétés importantes :
	\begin{itemize}
		\item \( e^0 = 1 \)
		\item \( e^{x+y} = e^x \cdot e^y \)
		\item \( e^{-x} = \frac{1}{e^x} \)
	\end{itemize}
	
	\subsection{Dérivée et primitive}
	
	La dérivée et la primitive de la fonction exponentielle sont données par :
	\[
	\frac{d}{dx} e^x = e^x \quad \text{et} \quad \int e^x \, dx = e^x + C.
	\]
	
	\subsection{Applications}
	
	La fonction exponentielle intervient dans de nombreux domaines :
	\begin{itemize}
		\item Croissance exponentielle : Modélisation de la croissance démographique, des intérêts composés, etc.
		\item Décroissance exponentielle : Modélisation de la désintégration radioactive, du refroidissement, etc.
	\end{itemize}
	                                                                    \section{Introduction aux fonctions logarithmiques}
	                                                                    
	                                                                    \subsection{Définition de la fonction logarithmique}
	                                                                    
	                                                                    La fonction logarithmique, notée \( \ln(x) \), est l'inverse de la fonction exponentielle. Autrement dit, pour tout réel positif \( y \), le logarithme naturel de \( y \) est le nombre \( x \) tel que :
	                                                                    \[
	                                                                    e^x = y \quad \text{soit} \quad x = \ln(y).
	                                                                    \]
	                                                                    
	                                                                    \subsection{Propriétés fondamentales}
	                                                                    
	                                                                    La fonction logarithmique possède plusieurs propriétés intéressantes :
	                                                                    \begin{itemize}
	                                                                    	\item \( \ln(1) = 0 \)
	                                                                    	\item \( \ln(ab) = \ln(a) + \ln(b) \) pour tous \( a > 0 \) et \( b > 0 \)
	                                                                    	\item \( \ln\left(\frac{a}{b}\right) = \ln(a) - \ln(b) \)
	                                                                    	\item \( \ln(a^b) = b \cdot \ln(a) \)
	                                                                    \end{itemize}
	                                                                    
	                                                                    \subsection{Dérivée et primitive de la fonction logarithmique}
	                                                                    
	                                                                    La dérivée de la fonction logarithmique est donnée par :
	                                                                    \[
	                                                                    \frac{d}{dx} \ln(x) = \frac{1}{x}, \quad \text{pour} \quad x > 0.
	                                                                    \]
	                                                                    Sa primitive est :
	                                                                    \[
	                                                                    \int \ln(x) \, dx = x \ln(x) - x + C.
	                                                                    \]
	                                                                    
	                                                                    \subsection{Exemples d'utilisation}
	                                                                    
	                                                                    \begin{itemize}
	                                                                    	\item Résolution d'équations exponentielles : Le logarithme permet de résoudre des équations comme \( e^x = a \), ce qui donne \( x = \ln(a) \).
	                                                                    	\item Applications : Le logarithme naturel est utilisé dans des domaines tels que la finance pour calculer des taux de croissance, ou en physique pour modéliser des phénomènes comme la décroissance radioactive.
	                                                                    \end{itemize}
	                                                                                                                                                                                  
	
	\section{Lien entre logarithme et exponentielle}
	
	\subsection{Relation fondamentale}
	
	La fonction exponentielle et la fonction logarithmique sont réciproques l'une de l'autre. Pour tout réel \( x \) et pour tout \( y > 0 \), on a :
	\[
	e^{\ln(x)} = x \quad \text{et} \quad \ln(e^x) = x.
	\]
	
	\subsection{Comportement asymptotique}
	
	En étudiant les limites des fonctions logarithmique et exponentielle, on obtient les résultats suivants :
	\[
	\lim_{x \to +\infty} e^x = +\infty, \quad \lim_{x \to -\infty} e^x = 0,
	\]
	\[
	\lim_{x \to 0^+} \ln(x) = -\infty, \quad \lim_{x \to +\infty} \ln(x) = +\infty.
	\]
	
	\subsection{Inversibilité}
	
	La fonction exponentielle étant strictement croissante et toujours positive, elle est bijective sur \( \mathbb{R} \). De même, la fonction logarithmique est strictement croissante et définie pour \( x > 0 \). Par conséquent, la fonction logarithmique est l'inverse de la fonction exponentielle, et réciproquement.
	
	
	
	
	
	
\end{document}